\documentclass[11pt]{article}
\usepackage[margin=1in]{geometry}
\usepackage{amsmath,amssymb,amsthm,booktabs,hyperref}
\hypersetup{hidelinks}
\newtheorem{theorem}{Theorem}
\newtheorem{lemma}[theorem]{Lemma}
\newtheorem{corollary}[theorem]{Corollary}
\newcommand{\E}{\mathbb{E}}
\newcommand{\Prb}{\mathbb{P}}
\newcommand{\pos}[1]{\left(#1\right)_{+}}
\title{Geometric Target Repair: Proven Bounds and the Remaining Depth Problem}
\author{Balance Trees project}
\date{November 2026}
\begin{document}
\maketitle
\begin{abstract}
We analyze the geometric-target local-repair implementation in this repository.
Each successful insertion preserves binary-search-tree order and takes at most two
rotations. The implementation draws fewer than two fair bits in expectation and
stores no per-node balance metadata. Its expected ancestor-walk cost is controlled
precisely by the insertion depth's excess over $\lceil\log_2 n\rceil$. Neither a
constant bound on that excess nor logarithmic expected depth holds against
adaptive insertion sequences, although both hold for monotone sequences.
We give a distribution-preserving sampler that is optimal in expected fair bits,
and a general logarithmic lower bound for some node's expected depth.
\end{abstract}

\section{Model and algorithm}
Keys are distinct and totally ordered; a duplicate is rejected before policy
state or random bits change. Start with an empty binary search tree (BST).
Depth counts edges from the root. A successful insertion first attaches the new
key as a leaf $v$, at depth $d$, and sets $n$ to the number of keys including
$v$. Write
\[ k=\lceil\log_2 n\rceil,\qquad L=k+1. \]
The strategy has a global successful-insertion counter; each node contains only
its key and structural child and parent links. The count needs $\Theta(\log n)$
bits (one word in the usual word-RAM model). We count comparisons for the BST
search separately. A rotation and a parent-pointer step have constant cost; a
fair-bit request costs one. Exact $k$ is obtained in constant word operations,
as with the implementation's integer bit length. Random-generator state is
external to tree nodes and is not free storage in an implementation.

Draw independent fair bits until the first zero or $L$ bits have been read. If
$m$ ones were read before stopping, set the target $T=\max(0,k-m)$. If $d\leq T$,
do nothing; if $d=T+1$, rotate $v$ up once. Otherwise walk from $v$ to its
ancestor $x$ at depth $T+2$. Let $p$ and $g$ be the parent and grandparent of
$x$. If $g,p,x$ form a straight path, rotate $p$ up once. If the path bends,
rotate $x$ up twice. All rotations are ordinary BST rotations. This is the
repository's \texttt{geometric\_target\_local\_repair} and
\texttt{target\_repair} code; it is distinct from the geometric stopping rules
studied by Brodal~\cite{brodal}.

Each bit in the mathematical model is independent of the preceding history and
of the current key, conditional on that history. We permit the next key to be
chosen adaptively from the observed past, before the new bits are drawn. This
covers fixed sequences. Python's seeded pseudorandom generator is deterministic;
the probability statements below describe the ideal fair-bit model, not a
cryptographic or information-theoretic guarantee about its finite seed.

\section{Deterministic properties}
\begin{theorem}\label{thm:structural}
After every successful insertion the result is a BST containing exactly the
previous keys and the new key, and at most two rotations have occurred. No
per-node balancing information is written. The claim holds for every bit stream
and for both fixed and adaptive insertion sequences.
\end{theorem}
\begin{proof}
Ordinary leaf insertion adds the new key in its unique search position. A single
rotation preserves the inorder sequence: for example, rotating a left child $p$
above its parent $g$ replaces the local order $A,p,B,g,C$ by the same order;
the other orientation is symmetric. Rotations neither add nor remove keys and
preserve parent/child connectivity. The depth cases use zero, one, one, or two
rotations respectively. In the last case $x$ has a parent and grandparent because
$T\geq0$ and $x$ is at depth $T+2$; in the bent case the first promotion leaves
$x$ directly below $g$, so the second is defined. The only persistent policy
state is the global count, and rotations update structural links only.
\end{proof}

\begin{lemma}[Exact visit count]\label{lem:visits}
Let $V$ be the number of \emph{counted ancestor visits} in the implementation's
repair routine, including the final inspected ancestor. Then, pointwise,
\[V=\pos{d-T-1}.\]
\end{lemma}
\begin{proof}
If $d\leq T+1$ the routine records no visits, and the right-hand side is zero.
Otherwise it traverses $d-T-2$ parent links and counts each departure, then
counts the selected ancestor once. Thus it records $d-T-1$ visits. Other constant
work, including the single-rotation case, is not counted as an ancestor visit.
\end{proof}

\section{Exact random-bit and repair costs}
\begin{theorem}[Target distribution and bits]\label{thm:bits}
For $k\geq0$, the target and number $B$ of requested bits satisfy
\begin{align*}
\Prb(T=0)&=2^{-k}, & \Prb(T=t)&=2^{-(k-t+1)}\quad(1\leq t\leq k),\\
\E B&=2-2^{-k}<2, & \E(k-T)&=1-2^{-k}.
\end{align*}
The statements hold conditional on any prior history and current key.
\end{theorem}
\begin{proof}
For $1\leq t\leq k$, the first $k-t$ bits must be one and the next zero.
Target zero occurs precisely when the first $k$ bits are all one, irrespective
of the following bit; for $k=0$ this event is certain. The $i$th bit is requested
exactly when the preceding $i-1$ bits are all one, for $1\leq i\leq k+1$.
Summing these probabilities gives $\E B=\sum_{j=0}^{k}2^{-j}=2-2^{-k}$.
Finally $k-T=\min(G,k)$, where $G$ is the number of leading ones in an
unbounded stream, so the tail-sum formula gives
$\E(k-T)=\sum_{j=1}^{k}2^{-j}=1-2^{-k}$. Conditional independence makes the
same calculation valid after every history.
\end{proof}

\begin{theorem}[Exact and bounded expected walking]\label{thm:walk}
Let $D$ be the new leaf's depth before repair. Conditional on a history and
current key, $D=d$ is fixed and the newly sampled $T$ is independent of it. If
$d\geq k+1$, then
\[\E[V\mid D=d]=d-k-2^{-k}.\]
For random $D$, including an adaptively chosen next key,
\[\E\pos{D-k-1}\ \leq\ \E V\ \leq\
\E\pos{D-k-1}+1-2^{-k}.\]
Consequently a uniform $O(1)$ expected repair-time bound is equivalent to a
uniform $O(1)$ bound on $\E\pos{D-k-1}$, under the stated cost model.
\end{theorem}
\begin{proof}
For $d\geq k+1$, $d-T-1\geq0$ for every target, so Lemma~\ref{lem:visits}
and Theorem~\ref{thm:bits} give
$\E V=d-k-1+\E(k-T)=d-k-2^{-k}$. For any $d$ and $T\leq k$,
\[\pos{d-k-1}\leq\pos{d-T-1}
\leq\pos{d-k-1}+(k-T).\]
Take conditional expectations and then average over the prior history. Sampling
and constant rotation work have expected $O(1)$ cost, while the parent loop
costs $\Theta(V)$ up to additive constants. The two inequalities therefore
prove the equivalence. An $O(\log n)$ bound on $\E D$ alone does not imply the
required additive-excess bound.
\end{proof}

\begin{theorem}[A reachable conditional-time counterexample]\label{thm:chain}
There is no uniform constant bound on expected repair time conditional on every
reachable history, even for a fixed insertion sequence.
\end{theorem}
\begin{proof}
Insert $1,2,\ldots,m$ and choose target zero at each insertion. Each such
choice has positive probability by Theorem~\ref{thm:bits}, so this finite
history has positive probability. Inductively the current maximum is the root:
the next key attaches as its right child and rotates above it. The resulting
tree is the descending left chain $m,m-1,\ldots,1$. Now insert zero. Its depth
is $m$, and every target is at most $k=\lceil\log_2(m+1)\rceil$. Therefore
Lemma~\ref{lem:visits} gives
$V\geq\pos{m-k-1}$ for every fresh target, proving the claim. The fixed
sequence is $(1,2,\ldots,m,0)$. This rare conditional history does not prove
that the \emph{unconditional} expected cost of that sequence, or any other
fixed sequence, is superconstant.
\end{proof}

\begin{theorem}[Monotone insertion sequences]\label{thm:monotone}
For strictly increasing insertions, the maximum key has depth at most
$\lceil\log_2 n\rceil$ after every repair, for every sequence of bits. Each
insertion uses at most one rotation, and its expected number of counted repair
visits, conditional on the past, is at most $1-2^{-k}<1$. The symmetric claims
hold for strictly decreasing insertions and the minimum key. Duplicate attempts
do not affect these bounds.
\end{theorem}
\begin{proof}
The root insertion has depth zero and needs no rotation. For $n\geq2$, assume
the previous maximum has depth at most $\lceil\log_2(n-1)\rceil\leq k$.
The new maximum attaches as its right child, so its insertion depth $d$ is at
most $k+1$. Its entire search path takes right edges. If $d\leq T$, there is no
rotation and $d\leq k$. Otherwise the single-leaf case rotates once, and the
two-edge case is straight, so it also rotates once. In either case, the rotated
child lies on the new maximum's path, so that maximum moves up exactly one
level; its final depth is at most $k$. This proves the induction and rotation
bound for every choice of target. Moreover $d\leq k+1$ implies
$V=\pos{d-T-1}\leq k-T$, whose conditional expectation is $1-2^{-k}$ by
Theorem~\ref{thm:bits}. Mirror left and right for decreasing insertions. This
result bounds the current extremal node.
\end{proof}

\begin{theorem}[Every node under monotone insertions]\label{thm:monotone-all}
For strictly increasing insertions, fix a key $x$ inserted at time $m$. After
$n\geq m$ insertions its expected depth satisfies
\[
\E[\operatorname{depth}_n(x)]\leq
2\lceil\log_2 m\rceil+\sum_{j=m+1}^{n}\frac1j=O(\log(n+1)).
\]
The same result holds for strictly decreasing insertions by symmetry.
\end{theorem}
\begin{proof}
The \emph{right spine} is the path obtained by following right children from
the root. Before each insertion let $a$ be the depth of $x$. If $x$ is on the
right spine, set $b=a$; otherwise let $b$ be the depth of its deepest right-spine
ancestor, whose left subtree contains $x$. By Theorem~\ref{thm:monotone}, the
old maximum is at depth at most $k$ for the next insertion's $k$, hence
$0\leq b\leq k$.

The new maximum is attached at depth $d$ on the right spine, with $d>b$.
Whenever the sampled target $t<d$, repair rotates the right-spine child at
depth $t+1$ above its parent at depth $t$. For the potential $a+b$ there are
three cases. If $t=b$, $x$ moves down one level and $b$ is unchanged. If
$t<b$, the branch point moves up one level, so $b$ decreases by one, and the
depth $a$ does not increase. In the boundary case where the branch point is
the rotated child, its left subtree may move across the rotation and retain
its depth; this still gives the stated inequality. If $t>b$, the rotation is
strictly below the relevant branch point and neither $a$ nor $b$ changes.
If no rotation occurs then $t\geq d>b$, so this last case still applies.

For $1\leq b\leq k$, Theorem~\ref{thm:bits} and a geometric sum yield
\[\Prb(T=b)=\Prb(T<b)=2^{-(k-b+1)}.\]
Consequently the conditional expected change of $a+b$ is at most zero.
For $b=0$ it is at most $\Prb(T=0)=2^{-k}\leq1/j$ at insertion $j$.
Immediately after $x$ is inserted, it is the maximum and lies on the spine;
Theorem~\ref{thm:monotone} gives $a=b\leq\lceil\log_2 m\rceil$. Summing the
conditional drift bounds over $j=m+1,\ldots,n$, and using $a\leq a+b$,
proves the result. The harmonic sum is at most $1+\ln n$, establishing the
asymptotic bound. Mirror left and right for decreasing insertions.
\end{proof}

The following adaptive construction uses the strong law of large numbers and
Fatou's lemma; see Williams~\cite{williams} for both results.
\section{An adaptive counterexample}

The following lower bound is for the algorithm specified above and the ideal
independent fair-bit model. It does not concern the distinct algorithms
analyzed by Brodal. The adversary sees the current tree before choosing each
next key, and all choices are made before the next insertion's random bits.

\begin{theorem}[Linear adaptive depth and repair]\label{thm:adaptive-lower}
There is an adaptive sequence of distinct rational keys and a key $x$ in the
initial prefix such that, after $n$ successful insertions,
$\mathbb E[\operatorname{depth}_n(x)]=\Omega(n)$. Along an infinite subsequence
of insertions, the expected number of counted repair visits is also $\Omega(n)$.
The expectations include all random bits, including those in the initial prefix.
\end{theorem}

\begin{proof}
Throughout, $x$ will be the largest inserted key, so its root-to-node path
consists entirely of right edges. Let $n$ include the next key, and write
$k=\lceil\log_2 n\rceil$. Assume for the moment that $k\geq2$ and that the
present depth $d$ of $x$ satisfies $d>k+2$. We specify two types of insertion.
All required gaps exist because the rational order is dense.
For definiteness, choose the midpoint of each bounded gap and one less than
the upper endpoint for an unbounded-below gap. The attack depends only on the
current tree, including its keys; the fixed schedule below is determined by
the number of nodes, which can be counted from that tree.

For an \emph{attack}, let $g$ be the node at depth $k$ on the path to $x$,
and insert a fresh key in $g$'s left subtree, smaller than every key already
in that subtree. If the subtree is empty, insert at its root gap. The insertion
path is $k$ right edges followed by one or more left edges. There are three
cases. With probability $1/2$, $T=k$: a left child of $g$ is promoted once,
so $g$, its right subtree, and $x$ move down one level. This includes the
case where the new leaf is the left child of $g$. With probability $1/4$,
$T=k-1$: the selected two-edge path is right--left, so the double rotation
leaves $g$ at its old depth and leaves $d$ unchanged. With probability $1/4$,
$T\leq k-2$: the selected path is right--right and the single rotation moves
$x$ up one level. Thus an attack changes $d$ by $+1,0,-1$ with respective
probabilities $1/2,1/4,1/4$.

For a \emph{probe}, insert a fresh key in the gap immediately below $x$,
choosing it within $x$'s left subtree. The insertion path reaches $x$, and
its insertion depth is at least $d+1$. Since $d>k+2$, every target $T\leq k$
selects a right--right path strictly above $x$. Exactly one rotation occurs,
moving $x$ up one level. By the exact visit formula, the probe has
\[ V=(d_{\mathrm{ins}}-T-1)_+\geq d-k. \]
Neither kind of key exceeds $x$, which remains the maximum.

Fix $m$ in advance. To initialize, insert $m$ distinct keys in decreasing
order, starting with $0$ and then choosing one less than the current minimum.
Afterward, the current maximum is $x$. At each subsequent step, count the
nodes to determine the next size and its position in a fixed schedule of nine
attacks followed by one probe. If the current depth of $x$ is at most $k+2$,
insert one less than the current minimum; otherwise use the scheduled attack
or probe. This is a deterministic function of the current tree and fixed $m$:
it does not remember whether the condition failed previously, inspect priorities
or random bits, or know future outcomes. On the event that
each sampled target in this prefix is zero, every new minimum rotates to the
root, leaving a right chain with $d_m(x)=m-1$. This event has strictly positive
probability for every fixed $m$.

It remains to choose a fixed $m$ so that the subsequent condition holds forever
with positive probability. Couple each future attack to an independent random
variable taking values $+1,0,-1$ with probabilities $1/2,1/4,1/4$; assign
$-1$ to each future probe. The coupling, used only in the analysis and not
available to the adversary, can read the first two latent fair bits
for an attack: zero gives $+1$, one then zero gives $0$, and two ones give
$-1$. Latent bits are independent even when a sampler stops before reading
them. Let $Z_s$ be the sum of the first $s$ such scheduled increments, with
$Z_0=0$. The sums of successive ten-increment blocks are independent and have
mean $9/4-1=5/4$. The strong law of large numbers, and the bounded remainder
within each block, imply
\[ \frac{Z_s}{s}\longrightarrow\frac18 \quad\text{almost surely}. \]
Consequently
\[ M=\sup_{s\geq1}\bigl(\lceil\log_2(s+1)\rceil-Z_{s-1}\bigr) \]
is finite almost surely. Choose a finite integer $B$ with $\Prb(M\leq B)>0$,
then choose a fixed $m\geq4$ large enough that
$m-1-\lceil\log_2m\rceil\geq B+3$. The initial-prefix event and $M\leq B$
are independent and their intersection $E$ has positive probability.

For step $s$ after the prefix, ceiling subadditivity and $m+s\leq m(s+1)$ give
\[\lceil\log_2(m+s)\rceil\leq
  \lceil\log_2m\rceil+\lceil\log_2(s+1)\rceil.\]
Inductively on $E$, before that step the depth is $m-1+Z_{s-1}$, so its
excess over the current $k$ is at least $3$. Hence the stated attack and probe
transitions apply at every step, and on $E$,
$d_{m+s}(x)=m-1+Z_s$ and $d_{m+s}(x)/(m+s)\to1/8$ almost surely.
Fatou's lemma yields
\[ \liminf_{n\to\infty}\frac{\mathbb E[d_n(x)]}{n}
   \geq \frac{\Prb(E)}8>0.\]
On the probe subsequence $n=m+10j$, the bound $V\geq d-k$, applied to the
depth just before the probe, and the same argument give
\[ \liminf_{j\to\infty}\frac{\mathbb E[V_{m+10j}]}{m+10j}
   \geq \frac{\Prb(E)}8>0.\]
Thus both stated lower bounds are unconditional.
\end{proof}

\begin{corollary}[Tight adaptive worst case]\label{cor:adaptive-tight}
Under the model of Theorem~\ref{thm:adaptive-lower}, the largest expected
node depth attainable by an adaptive adversary is $\Theta(n)$. The expected
repair visits in that theorem are $\Theta(n)$ along its probe subsequence.
\end{corollary}
\begin{proof}
Every node of every $n$-node binary tree has depth at most $n-1$, and the
counted visits for an insertion are at most its insertion depth. These
deterministic upper bounds, together with Theorem~\ref{thm:adaptive-lower},
give the claimed matching bounds.
\end{proof}

The proof uses an infinite dense key order. The Python implementation compares
keys at runtime, although its annotations say \texttt{int}; the rational model
must therefore be distinguished from the annotated interface. For any fixed
finite horizon, its finitely many possible adaptive histories use finitely many
rational keys; an order-preserving rescaling embeds those keys in the integers.
This yields the same finite-horizon worst-case bounds over integer inputs, with
the spacing allowed to depend on the horizon. Comparisons here have unit cost;
the integer spacing may require more than $O(\log n)$ bits per key, so the
embedding does not establish the result for a fixed-width integer universe or
account for arbitrary-precision arithmetic costs. The counterexample does not
prove a linear lower bound for a fixed, nonadaptive insertion sequence.


\section{An entropy-optimal alternative sampler}
The following alternative changes only bit consumption, not the target law:
read bits until the first zero or until $k$ bits have been read, then return
$T=k-m$ where $m$ is the number of leading ones read. When $k=0$, return zero
without reading a bit. The original sampler requests one more bit whenever its
first $k$ bits are ones, although the target is already known.

\begin{theorem}[Optimal expected bits for this target law]\label{thm:entropy}
The alternative has exactly the target distribution in Theorem~\ref{thm:bits}
and uses $2-2^{1-k}$ expected bits. It attains the entropy $H(T)$, hence no
exact fair-bit prefix sampler for one independent target can use fewer bits in
expectation without recycling randomness from previous or future samples.
\end{theorem}
\begin{proof}
The alternative stops at the prefix $1^{k-t}0$ for $1\leq t\leq k$ and at
$1^k$ for $t=0$. These prefixes have the probabilities stated in
Theorem~\ref{thm:bits}, and their lengths are $-\log_2\Prb(T=t)$. The expected
length is $\sum_{j=0}^{k-1}2^{-j}=2-2^{1-k}$, also valid as an empty sum for
$k=0$; it equals $H(T)$.

For the lower bound, let $S$ be the random stopping prefix of any exact
almost-surely terminating fair-bit sampler. Its possible prefixes are
prefix-free, and a prefix $s$ of length $|s|$ occurs with probability
$2^{-|s|}$. Hence $H(S)=\E|S|$ (also in the extended sense). Since the output
is a deterministic function of $S$, $H(T)\leq H(S)$ by entropy's chain rule;
see Shannon~\cite{shannon}. Thus the alternative meets the bound. This is
optimal for the \emph{given output distribution}, not a lower bound on all
possible balancing algorithms. Bit recycling across outputs changes the model.
\end{proof}
The alternative is a mathematical variant, not the implemented sampler. Even
with the same pseudorandom seed it generally changes subsequent draws and tree
shapes. The repository's stored-priority treap reference counts 32 logical bits
per successful insertion and stores a priority per key; it is not the same model
as ideal continuous-priority treaps~\cite{aragon}.

\section{Depth lower bound and research status}
\begin{theorem}[Universal logarithmic lower bound]\label{thm:lower}
For any distribution on binary trees containing the same $n\geq4$ distinct keys,
there exists a key whose expected depth is at least
$\frac12(\lfloor\log_2 n\rfloor-1)$. Thus a uniform bound on all nodes' expected
depth cannot be $o(\log n)$, even with arbitrary balancing metadata and rotations.
\end{theorem}
\begin{proof}
Put $h=\lfloor\log_2 n\rfloor-1$. Every binary tree has at most $2^h-1<n/2$
vertices at depths below $h$. Hence at least $n/2$ vertices have depth at least
$h$, so the average of all $n$ depths is at least $h/2$ in every realization.
Take expectations and interchange the finite sum with expectation. At least one
of the $n$ fixed keys has expected depth at least the average.
\end{proof}

\begin{center}
\begin{tabular}{@{}lp{0.55\linewidth}@{}}
\toprule
Property & Status for the implemented strategy \\
\midrule
BST and key preservation & Proved, Theorem~\ref{thm:structural}. \\
Node balance metadata & None; a global count and external random state are used. \\
Rotations & At most two per successful insertion, deterministically. \\
Random bits & Fewer than two in expectation, Theorem~\ref{thm:bits}. \\
Locality & Uses $n$ and may traverse far above the inserted node. \\
Expected repair time & Adaptive sequences can require $\Omega(n)$ expected visits by Theorem~\ref{thm:adaptive-lower}; the fixed-sequence bound is open. \\
Expected depth of every node & $O(\log n)$ for monotone sequences by Theorem~\ref{thm:monotone-all}, but $\Omega(n)$ is possible under adaptive keys by Theorem~\ref{thm:adaptive-lower}. The general fixed-sequence case is open. \\
Optimality & The logarithmic lower bound is achieved for monotone sequences but not against adaptive insertions. \\
\bottomrule
\end{tabular}
\end{center}
The guarantees in the first four rows hold for fixed and adaptive histories under
the stated fresh-bit assumption. The counterexample directly analyzes the
target-depth rule; Brodal~\cite{brodal} exhibits related linear-depth behavior
for different bottom-up random repair rules. Constant expected repair needs
the excess-depth estimate in Theorem~\ref{thm:walk}. Finite benchmarks and
bounded enumeration can suggest conjectures or refute finite universal claims,
but do not establish asymptotic bounds.

\begin{thebibliography}{9}
\bibitem{brodal} Gerth St\o lting Brodal, ``Bottom-up Rebalancing Binary Search
Trees by Flipping a Coin,'' 2024. \url{https://arxiv.org/abs/2404.08287}.
\bibitem{aragon} Cecilia R. Aragon and Raimund G. Seidel, ``Randomized Search
Trees,'' \emph{FOCS}, 1989. \url{https://faculty.washington.edu/aragon/pubs/rst89.pdf}.
\bibitem{shannon} Claude E. Shannon, ``A Mathematical Theory of Communication,''
\emph{Bell System Technical Journal} 27, 1948.
\url{https://people.math.harvard.edu/~ctm/home/text/others/shannon/entropy/entropy.pdf}.
\bibitem{williams} David Williams, \emph{Probability with Martingales},
Cambridge University Press, 1991. \url{https://doi.org/10.1017/CBO9780511813658.028}.
\end{thebibliography}
\end{document}
